\section{Resultados y discusión}
\label{sec:resultados}

Ambos niveles se resolvieron para los 30 polígonos (183.01 ha) con $\alpha=0.05$, $\theta=0.05$, $\delta=0.10$ y $K=10{,}000$ escenarios; el mosaico es de $4\times13=52$ posiciones al tresbolillo y la matriz es la de la \cref{tab:matriz-competencia}. Todos los resultados se reproducen con la herramienta \texttt{refor} del repositorio (semilla 42).

\subsection{Compra}
\label{subsec:res-compra}

Con $B=\$12{,}500$/ha ningún polígono es factible. El presupuesto mínimo que satisface \eqref{eq:total}--\eqref{eq:probabilidad} es, en promedio ponderado por superficie, $B^*=\$14{,}496$/ha (entre \$14{,}361 y \$14{,}934), 16\% más que el disponible. La infactibilidad es estructural y no se resuelve aceptando más riesgo: $B^*$ apenas baja de \$14{,}664/ha con $\alpha=0.01$ a \$14{,}117/ha con $\alpha=0.50$, un ahorro menor al 4\% a cambio de multiplicar por cincuenta la probabilidad de incumplimiento. Lo que excede el presupuesto es el déficit mismo, cuyo costo se estimó en la \cref{subsec:compra-formulacion}, no el margen de seguridad. Los polígonos pequeños son los más caros por hectárea (\$14{,}934 en el polígono 1, de 1.28 ha), porque el margen necesario para cumplir con probabilidad $1-\alpha$ crece más lento que la meta y pesa más cuando ésta es chica.

La compra sin restricción de presupuesto suma 108{,}680 plantas y \$2.65 millones (594 plantas/ha). Junto con las 129 plantas/ha existentes, la densidad al plantar es de 723 plantas/ha, 10\% sobre la meta: es la sobreplantación que absorbe la mortalidad del primer año con la confiabilidad exigida. El costo esperado de reposición es de solo \$5/ha, de modo que con $\alpha=0.05$ el costo es casi todo de adquisición. La composición óptima agota la banda \eqref{eq:banda} en la dirección que anticipa el costo por planta sobreviviente $c_s/q_s$: las cuatro especies de \textit{Opuntia} suben 10\%, \textit{P.~laevigata}, \textit{Y.~filifera} y \textit{A.~scabra} bajan 10\%, \textit{A.~lechuguilla} y \textit{A.~striata} 8\%, y \textit{A.~salmiana} se mantiene.

\subsection{Distribución espacial}
\label{subsec:res-espacial}

Al redondear la composición de cada polígono a 52 posiciones, los 30 polígonos se reducen a seis mosaicos distintos; el primero lo comparten 18 polígonos (\cref{tab:res-mosaicos}). El algoritmo memético reduce la competencia entre 14.5\% y 16.3\% respecto a una asignación aleatoria de igual composición. El resultado es muy estable: la desviación estándar de las 30 corridas no excede 0.006 y la mejora se detiene, en promedio, antes de la generación 25, por lo que el criterio de paro de 100 generaciones sin mejora es holgado.

Las claves obtenidas eliminan toda vecindad intraespecífica salvo una o dos entre plantas de \textit{A.~salmiana}, que ocupa 27--29\% de las posiciones, de 156 pares vecinos por mosaico. \textit{P.~laevigata}, con la mayor penalización intraespecífica (1.50), nunca queda junto a otra de su especie y actúa como separador: cerca de la mitad de sus vecinas son \textit{A.~salmiana} y un tercio \textit{O.~streptacantha}, las dos especies abundantes con penalización intraespecífica alta, con las que la suya es mínima (0.20 y 0.17).

\begin{table}[t]
\centering
\caption{Mosaicos de 52 posiciones. $n_p$: polígonos que comparten el mosaico; memético: mejor competencia de 30 corridas (media $\pm$ d.e.); coef.: cociente respecto al azar; cota: cota inferior del modelo entero en 300 s; brecha: distancia máxima del memético al óptimo.}
\label{tab:res-mosaicos}
\small
\setlength{\tabcolsep}{3pt}
\begin{tabular}{@{}crcccc@{}}
\toprule
Mosaico & $n_p$ & Memético & Coef. & Cota & Brecha \\
\midrule
1 & 18 & 78.18 ($78.19\pm0.01$) & 0.838 & \textsc{pend} & \textsc{pend} \\
2 &  5 & 77.37 ($77.39\pm0.01$) & 0.838 & \textsc{pend} & \textsc{pend} \\
3 &  1 & 81.29 ($81.29\pm0.00$) & 0.855 & \textsc{pend} & \textsc{pend} \\
4 &  1 & 76.77 ($76.78\pm0.01$) & 0.837 & \textsc{pend} & \textsc{pend} \\
5 &  4 & 76.80 ($76.80\pm0.00$) & 0.837 & \textsc{pend} & \textsc{pend} \\
6 &  1 & 81.85 ($81.85\pm0.00$) & 0.853 & \textsc{pend} & \textsc{pend} \\
\bottomrule
\end{tabular}
\end{table}

% PENDIENTE: comparación con el modelo entero (\cref{subsec:espacial-entera}).
